\section{Функциональные требования}

\subsection{Загрузка и хранение документов}

\subsubsection{Выбор типа документа}
Пользователь при загрузке указывает тип документа: \emph{шаблон ВКР} или \emph{отчёт по практике}. Выбор возможен для одного файла или для пакетной загрузки (общий тип для группы файлов или индивидуальный, в зависимости от реализации).

\subsubsection{Загрузка файлов}
Поддерживается загрузка одного PDF-файла или нескольких файлов за один раз (drag-and-drop и выбор через диалог). Принимаются только файлы в формате PDF; при загрузке проверяется расширение и соответствие содержимого формату (magic bytes). Файлы с неподходящим форматом отклоняются с понятным сообщением об ошибке.

\subsubsection{Безопасность загружаемых PDF (обязательное требование)}
Перед приёмом и сохранением каждый загруженный файл должен проходить проверку безопасности. Не допускается приём файлов, содержащих вредоносное ПО, исполняемые объекты или типичные PDF-эксплойты. Проверки включают:
\begin{itemize}
  \item валидацию структуры PDF и отклонение некорректных или подозрительных объектов (встроенные файлы, JavaScript, действия запуска, внешние ссылки/URI, которые могут быть использованы во вред);
  \item при необходимости~--- интеграцию с антивирусной проверкой (например, ClamAV или корпоративное решение) в соответствии с политикой ИТМО.
\end{itemize}
Файлы, не прошедшие проверку, не сохраняются; пользователю возвращается однозначное сообщение о причине отклонения (безопасность).

\subsubsection{Ограничение размера и хранение}
Устанавливается настраиваемый максимальный размер файла (по умолчанию~--- 50\,МБ). Сохраняются только файлы, прошедшие проверку формата и безопасности. Хранение файлов выполняется в S3-совместимом хранилище; метаданные и результаты проверки~--- в базе данных. Реализуется просмотр истории загрузок и списка проверенных документов с указанием типа и статуса.

\subsection{Правила проверки (два типа документов)}

\subsubsection{Проверка шаблона ВКР (PDF)}

Формальная проверка выполняется в соответствии с документом «Требования к выпускным квалификационным работам» Университета ИТМО. Оцениваются в том числе:

\begin{itemize}
  \item \textbf{Структура:} наличие и порядок элементов (титульный лист, задание, аннотация, содержание; при необходимости~--- списки сокращений и терминов; введение, основная часть, заключение; список использованных источников; приложения).
  \item \textbf{Общее оформление:} формат страницы А4, поля (левое 30\,мм, правое 15\,мм, верхнее и нижнее 20\,мм), абзацный отступ 1,25\,см, шрифт Times New Roman не менее 12\,пт (рекомендуется 14\,пт), межстрочный интервал 1,5, выравнивание по ширине.
  \item \textbf{Нумерация страниц:} арабские цифры, сквозная, внизу по центру.
  \item \textbf{Заголовки:} структурные элементы (СОДЕРЖАНИЕ, ВВЕДЕНИЕ, ЗАКЛЮЧЕНИЕ и т.\,д.)~--- по центру, прописными, без номера; разделы и подразделы~--- нумерация и оформление по документу; каждый структурный элемент и глава начинаются с новой страницы.
  \item \textbf{Иллюстрации:} нумерация, подпись в формате «Рисунок N~--- Наименование», размещение.
  \item \textbf{Таблицы:} нумерация, подпись в формате «Таблица N~--- Наименование», размещение.
  \item \textbf{Формулы:} нумерация, размещение, пояснения под формулой.
  \item \textbf{Списки:} сокращений и условных обозначений; терминов; список использованных источников~--- оформление и наличие в содержании.
  \item \textbf{Приложения:} обозначение (А, Б, …), заголовки, перечисление в содержании.
\end{itemize}

Правила проверки настраиваются с опорой на эталонные материалы.

\subsubsection{Проверка отчёта по практике (PDF)}

Проверяется соответствие формальным требованиям к оформлению отчётных документов и структуре задания (подробное описание этапов 2, 3 и 4; в отчёте~--- результат только этапа 4; для этапа 2~--- наличие скриншотов присутствия). Оцениваются наличие описания этапов, наличие скриншотов, фокус основного текста на результате этапа 4. Глубокая смысловая оценка содержания не выполняется. Правила проверки настраиваются с опорой на эталонные материалы.

\subsection{Результаты проверки и отчётность}

После загрузки и проверки результаты отображаются в интерфейсе по каждому документу и по каждой проверке. Формируется детальный отчёт: перечень пройденных и непройденных проверок, предупреждения. Статус документа: «Пройдено» / «Не пройдено» / «Предупреждение» с цветовой индикацией. Реализуется экспорт отчёта о проверке (например, PDF или CSV). При пакетной загрузке выводится сводка по всем загруженным файлам.

\subsection{Отображение и отправка результатов проверки}

Реализуется панель (дашборд) со списком документов по типам (шаблон ВКР / отчёт по практике), статусу и дате. Предусматриваются фильтрация и поиск по типу документа, статусу, дате. Система автоматически формирует отчёт с результатами формальной проверки (предупреждения, ошибки) и предоставляет его пользователю (просмотр в интерфейсе, экспорт). Ручная проверка и комментарии преподавателя по загруженным отчётам и шаблонам не предусмотрены. В интерфейсе доступен просмотр PDF с подсветкой мест, связанных с замечаниями. Отображается статистика: доля успешно прошедших проверку, типичные ошибки.

\subsection{Интерфейс на основе электронной почты}

Помимо веб-интерфейса, предусмотрен интерфейс, использующий электронную почту. Система обрабатывает входящие письма с заданными темами (настраиваемый список или шаблоны), извлекает PDF-файлы из вложений, выполняет формальную проверку по выбранному или определённому типу документа и автоматически отправляет отправителю (или привязанному студенту) отчёт с результатами проверки (предупреждения, ошибки).

Возможные варианты реализации:
\begin{itemize}
  \item Опрос почтового ящика по протоколу IMAP (или аналог с OAuth2 для корпоративной почты ИТМО): периодическая проверка папки «Входящие», фильтрация писем по теме (темам), извлечение вложений с расширением PDF, запуск проверки, формирование отчёта и отправка ответным письмом на адрес отправителя.
  \item Приём уведомлений о новых письмах через webhook или очередь сообщений (если почтовый сервер или промежуточный сервис поддерживает такие механизмы): при появлении письма с заданной темой~--- извлечение вложений, проверка, отправка отчёта.
\end{itemize}
Тип документа (шаблон ВКР / отчёт по практике) может определяться по теме письма, по ключевым словам в теме или задаваться конфигурацией для используемого почтового ящика.

\subsection{Управление пользователями и аутентификация}

Аутентификация предпочтительно выполняется через \textbf{itmo.id}; при невозможности~--- по механизму, согласованному с инфраструктурой ИТМО. Реализуется разграничение прав (роли: студент; при необходимости~--- преподаватель, администратор). Профиль пользователя минимален (имя, email из itmo.id).

\subsection{Настройка}

Настраиваются правила проверки: включение/отключение или настройка отдельных проверок по типу документа (ВКР / практика). Через конфигурацию (в т.\,ч. pydantic-settings) задаются ограничение размера файла, срок хранения, параметры S3 и БД.
